% !PRINTED 408
\weberpage{408}
%<<CONT
rational functions of $x_0, x_1, \ldots, x_{n-1}$ which arise from one of
them by the permutations of the alternating group, and are roots of the
equation

\weq{1}{F\,(u, z) = 0}

where $z$ is a rational function of the $x$ which remains unchanged by the
permutations of the alternating group.

We prove the following second lemma:

\gesperrt{2. If the rational function $\Psi\,(u, u_1, \ldots, u_{\nu-1})$
vanishes identically when for the
\[
x_0, x_1, \ldots, x_{n-1}
\]
certain rational functions of a variable $t$}

\weq{2}{x_h = \varphi_h\,(t)}

\gesperrt{are put, and if by this substitution $z$ does not become
independent of $t$, then $\Psi$ vanishes identically also as a function of
the independent variables $x_0, x_1, \ldots, x_{\nu-1}$.}
% [sic] The source prints x_{\nu-1} here, although the independent variables were introduced through x_{n-1}.

The Galois group of equation (1) in the field of the rational functions of
$z$ will consist of certain permutations of the roots
$u, u_1, \ldots, u_{\nu-1}$. If we carry out these permutations in the
function $\Psi$, and form the product

\weq{3}{\Pi\,\Psi\,(u, u_1, \ldots, u_{\nu-1}) = f\,(z)}

of all the functions so obtained, there results a rational function
$f\,(z)$ of $z$. Now if one of the factors of the product (3) vanishes
identically after the substitution (2), and $z$ is not independent of $t$,
then $f\,(z)$ must be identically equal to zero, and consequently one of
the factors of the product (3) must also vanish identically as a function
of the $x$. But if one of these factors is identically equal to zero, then
all the other factors vanish as well, because in every rational equation
between the $u$ one can carry out all the permutations of the Galois
group. With this the theorem 2. is proved.

This theorem, taken together with the theorem 1. proved in the preceding
section, now gives the following result:

\gesperrt{3. If $u, u_1, \ldots, u_{\nu-1}$ are the roots of an equation
(1) depending on a parameter $z$, and if $u, u_1, \ldots,
u_{\nu-1}, z$ are rational functions of the}
%CONT>>
